\section{Algorithmic Composition and Structural Abstraction}\label{sec:algorithmic-composition}

Algorithmic composition is the technique of using algorithms---formal sets of rules or instructions---to create music.
While this concept has roots in classical music (such as Mozart's musical dice games), the advent of computers has
exponentially expanded its scope and capabilities~\cite{Nierhaus_2009}.

\subsection{Procedural Generation vs. Direct Input}\label{detail-subsec:procedural-generation-vs.-direct-input}

Traditional digital music production relies on direct input.
A composer uses a MIDI keyboard or a graphical piano roll in a digital audio workstation to manually input every
note, chord, and rest.
This process grants total control over the micro-details of a performance but can be tedious when constructing complex,
repetitive, or mathematically structured sections.

Algorithmic composition shifts the paradigm to procedural generation.
Instead of writing the notes, the composer writes the \textit{rules} that generate the notes.
For example, rather than manually copying and pasting a bassline across 32 bars while applying slight variations to the
velocity, a composer might write a simple loop that algorithmically modifies the velocity parameter on each iteration.

\subsection{Structured Abstraction and Source Compression}\label{detail-subsec:structured-abstraction}

The true power of algorithmic composition lies in its ability to abstract musical ideas.
A chord progression can be represented as an array of data points, and a rhythm can be defined as a mathematical
function over time.
This approach allows composers to experiment rapidly: changing a foundational variable (such as a root note or a time
signature) can instantly cascade throughout the entire algorithmic structure, generating a new variation of
the piece without the need for manual editing.

Patterns and loops provide a form of \textit{source-level compression}.
Consider a motif defined once and invoked inside a loop that runs 256 times.
The source program names the motif and the repetition count; the compiler expands both into a flat timeline of note
events at lowering time.
Listing~\ref{lst:compression-example} illustrates the idea with Motivo syntax: three lines of pattern definition plus
a loop call produce a much longer MIDI event stream.

\begin{lstlisting}[label={lst:compression-example},caption={A compact Motivo program whose lowering expands into many MIDI note events.}]
pattern motif(int dur) {
    play A4 :dur;
    play rest :dur;
    play C5 :dur;
}

track melody using piano {
    loop (256) { play motif(1); }
}
\end{lstlisting}

Each call to \texttt{motif(1)} contributes three played values (two notes and one rest).
After 256 iterations the lowering stage materializes 768 played values on a single track, each mapped to one or more
MIDI messages.
The composer edits the motif definition or the loop bound---not hundreds of individual piano-roll cells.
Motivo aims to provide a type-checked environment specifically optimized for this methodology of structured,
algorithmic creation.
